\begin{thesisindex}
% Optional manual index: sort entries yourself; pages remain live references.
\gzuindexitem{交叉引用}{chap:introduction}
\gzuindexitem{三线表}{tab:parameters}
\gzuindexitem{补充材料}{app:extra}
\end{thesisindex}

\begin{resume}
\section{教育经历}
20XX 年 9 月至 20XX 年 6 月，在某大学学习，获学士学位。

20XX 年 9 月至今，在贵州大学攻读硕士学位。

\section{攻读硕士学位期间发表的学术论文及其他成果}
请填写实际成果及其发表或录用状态。下列条目仅演示格式，不是实际论文。
\begin{achievements}
  \item 作者列表. 实际论文题名. 期刊名称, 年份, 卷(期): 页码.（请据实替换）
  \item 其他成果名称及说明.（无相应成果时删除）
\end{achievements}

\section{获奖情况}
请填写获奖时间、奖项名称和授奖单位；无相关内容时可删除本节。

\section{参加项目}
请填写项目名称、起止时间及承担的工作；避免将未参与的工作列入个人经历。
\end{resume}

\begin{postscript}
本页为可选后记示例。可记录论文写作过程中的补充说明，也可删除整个后记环境。后记不使用附录编号，其页码继续接续前面的内容。
\end{postscript}
